% ---------------------------------------------------------------------
\subsection{Headphones: the MagLS binaural decoder (\fn{Binaural.cpp})}
\label{sec:math-binaural}
The Ambisonic field $a_k(t)$, $k = 0\dots K-1$, $K=(L+1)^2$, is turned into the two ear signals by $2K$ filters,
$e(t) = \sum_k d_{e,k} * a_k$ for $e\in\{\text{left},\text{right}\}$. \fn{BinauralFilters::design} computes these
filters from a set of head-related impulse responses (FABIAN or a SOFA file) by the magnitude least-squares method
\citep{schoerkhuber2018magls}, the method of the IEM BinauralDecoder. It runs once on a background thread.

\paragraph{Directions and weights.} A set with more than 3000 directions is thinned to about 3000 by taking, for each
point of a Fibonacci lattice $\bigl(\arccos(1 - 2\frac{i+1/2}{n}),\ \pi(1+\sqrt5)(i+\frac12)\bigr)$, the nearest
measured direction (the lattice spreads points almost evenly). Measured grids are rarely uniform (dense at the
horizon, sparse below), so every direction $i$ gets a quadrature weight inversely proportional to the number of
directions in a spherical cap around it that would hold about 12 points of an even grid,
\begin{equation}
  w_i \propto \frac{1}{\#\{j : \mathbf u_i\cdot\mathbf u_j > 1 - 2\cdot\min(\tfrac12, 12/Q)\}},\qquad \sum_i w_i = 1,
\end{equation}
so a region measured twice as densely counts as much as a sparse one, not twice as much.

\paragraph{The least-squares projection.} With the matrix $\mathbf Y\in\mathbb R^{Q\times K}$ of the plugin's own
encoder at the measured directions ($Y_{ik} = Y_k(\varphi_i,\vartheta_i)$, so the decoder fits exactly what is
encoded) and $\mathbf W = \operatorname{diag}(w_i)$, the coefficients that best explain a function $t_i$ given on the
directions are the weighted least-squares solution
\begin{equation}
  \mathbf d = \bigl(\mathbf Y^{\mathsf T}\mathbf W\mathbf Y + \lambda\mathbf I\bigr)^{-1}\mathbf Y^{\mathsf T}\mathbf W\,\mathbf t,
  \qquad \lambda = 10^{-3}\,\operatorname{tr}(\mathbf Y^{\mathsf T}\mathbf W\mathbf Y)/K .
\end{equation}
The matrix $\mathbf P = (\mathbf Y^{\mathsf T}\mathbf W\mathbf Y+\lambda\mathbf I)^{-1}\mathbf Y^{\mathsf T}\mathbf W$ is
computed once by a Cholesky factorisation (the regularised normal matrix is symmetric positive definite) and then
applied at every frequency and ear; the small Tikhonov term $\lambda$ keeps it well conditioned where the grid has
holes.

\paragraph{The bulk delay.} The time until the sound reaches the nearer ear is the same for all directions and carries
no information. It is found as the earliest sample over all HRIRs where $|h|$ exceeds 10\,\% of that HRIR's peak,
minus two samples, $\tau_b$, and removed from every spectrum by $H(f)\,e^{+i2\pi f\tau_b}$ before the fit. It is put
back afterwards as $\tau_b + \SI{0.5}{ms}$, the half millisecond being headroom for the slight pre-ringing of the
fitted filters. Without this the phase above the cut-off (below) would put the high frequencies at $t=0$, before the
low ones, and the filter would wrap around its end.

\paragraph{Magnitude least squares.} Below the frequency $f_c$ at which the order $L$ can still describe the head,
\begin{equation}
  f_c = \operatorname{clamp}_{[500,\,4000]}\Bigl(\frac{L\,c}{2\pi r_h}\Bigr),\qquad r_h = \SI{8.75}{cm},
\end{equation}
($kr \le L$, the usual rule for the order needed for a sphere of radius $r$), the HRTFs are fitted as they are,
$\mathbf d(f) = \mathbf P\,\mathbf h(f)$. Above it a low order cannot follow the fast changes of phase with
direction, and a complex fit would cancel itself into a quiet, dull sound. MagLS fits only the magnitude there and lets
the phase be whatever is smoothest: the target at bin $b$ is
\begin{equation}
  t_i(f_b) = |h_i(f_b)|\cdot\frac{\hat h_i(f_{b-1})}{|\hat h_i(f_{b-1})|},\qquad
  \hat h_i(f) = \sum_k Y_{ik}\,d_k(f),
\end{equation}
the measured magnitude with the phase the fit itself had one bin lower. Going up bin by bin, the phase continues
smoothly from the last exact bin, which is what makes the solution converge. After the fit the coefficients of a bin
above $f_c$ are scaled by $g = \operatorname{clamp}_{[0.25,4]}\sqrt{E_T/E_D}$, the ratio of the diffuse-field energies
of the target and of the fit, $E = \sum_i w_i|h_i|^2$: a low-order fit loses energy where it smooths, and this puts it
back (the diffuse-field covariance constraint of the MagLS paper, in its simplest scalar form). DC and Nyquist are
made real.

\paragraph{Diffuse-field equalisation.} Optionally (\ui{Diffusfeld-Entzerrung}) the average of both ears' diffuse-field
energy is smoothed over a third of an octave (bins $b/1.1225\dots1.1225b$), compared with its mean between 200 and
\SI{1000}{Hz}, and the inverse, clamped to $\pm\SI{12}{dB}$ and held constant below \SI{150}{Hz}, is applied as a
\emph{minimum-phase} filter. The minimum-phase version of a magnitude $|E(f)|$ is obtained through the real cepstrum:
$c = \mathcal F^{-1}\{\ln|E|\}$, fold it onto positive time ($c_0$, $2c_n$ for $0<n<N/2$, $c_{N/2}$, zero after),
transform and exponentiate, $E_\mathrm{min} = \exp\bigl(\mathcal F\{\text{folded } c\}\bigr)$. A minimum-phase filter adds
no latency and no pre-ringing, which a linear-phase equaliser of this resolution would.

\paragraph{Level and length.} All filters are scaled so that the diffuse-field energy of each ear between 300 Hz and
\SI{3}{kHz}, $\mathbf d^{\mathsf H}(\mathbf Y^{\mathsf T}\mathbf W\mathbf Y)\mathbf d$ averaged over the band, is $\tfrac12$:
the same $-\SI{3}{dB}$ per ear as an equal-power stereo pan, so switching from stereo to headphones does not change
the loudness. The filters are cut to 512 samples at \SI{48}{kHz} (\SI{10.7}{ms}) with a raised-cosine fade over the
last eighth, and resampled when the host runs at another rate. At run time they are one lag-1 stage of
\fn{PartitionedConvolution} with blocks of 128 samples, $K$ inputs and two outputs.

\paragraph{Turning the head (\fn{AmbiYawRotator}).} Before decoding, the field is turned by the head's yaw $\psi$
from the 3D view (and, since 15 November, the AmbiX outputs too): turning the head to the left turns the scene to the
right, so the rotation of \S\ref{sec:math-ir} is applied with the angle $-\psi$. The angle is interpolated across the
block, $\psi(i) = \psi_0 + \operatorname{wrap}(\psi_1-\psi_0)\,(i+1)/N$, and the $\cos m\psi$, $\sin m\psi$ are built by the
angle-addition recurrence from one $\sin/\cos$ pair per sample; wrapping the difference makes a turn through
$\pm180^\circ$ take the short way.

% ---------------------------------------------------------------------
\subsection{Two loudspeakers: crosstalk cancellation (\fn{Transaural.h})}
\label{sec:math-transaural}
Over two loudspeakers at $\pm\alpha$ each ear hears both. With the near-ear and far-ear HRTFs $H_n$, $H_f$ of the
speaker direction and the ratio $R = H_f/H_n$, the symmetric path from the speakers $(s_L, s_R)$ to the ears is
\begin{equation}
  \begin{pmatrix} e_L \\ e_R \end{pmatrix} = H_n \begin{pmatrix} 1 & R \\ R & 1 \end{pmatrix}
  \begin{pmatrix} s_L \\ s_R \end{pmatrix}.
\end{equation}
The matrix has the eigenvectors $(1,1)$ and $(1,-1)$ with the eigenvalues $1+R$ and $1-R$, so its inverse splits into a
\emph{sum} and a \emph{difference} filter: with $m = (b_L+b_R)/2$ and $s = (b_L-b_R)/2$ of the binaural pair,
\begin{equation}
  s_L = C_s m + C_d s,\quad s_R = C_s m - C_d s,\qquad C_s \approx \frac{1}{1+R},\quad C_d \approx \frac{1}{1-R}.
\end{equation}
Two convolutions instead of four, and the common factor $1/H_n$ (the colour of the speakers' own direction) is left out
or treated separately (below).

\paragraph{The ratio $R$ (\fn{TransauralFilters::design}).} $R$ is measured from the HRTF set at $\pm\alpha$ (both
speakers, averaged). Its phase contains the interaural time difference $\tau$, found as the lag of the maximum of the
cross-correlation of the near-ear and far-ear impulse responses (within \SI{1.5}{ms}), refined to a fraction of a
sample by fitting a parabola through the maximum and its neighbours,
$\delta = \tfrac12(a-c)/(a-2b+c)$. The delay is taken out, $R\,e^{+i2\pi f\tau}$, the remainder smoothed over a twelfth
of an octave as a complex number (real and imaginary parts: sound bending round the head changes its delay with
frequency, and a pure delay alone cancels only about \SI{6}{dB}), limited to $|R|\le0.97$, and the delay put back.

\paragraph{Regularised inversion.} At low frequencies the head is small against the wavelength, $R\to1$ and
$1/(1-R)$ grows without limit. The inverse is therefore regularised,
\begin{equation}
  C_x = \frac{\overline{X}}{|X|^2 + \beta},\qquad X\in\{1+R,\ 1-R\},\qquad \beta = \frac{1}{4g^2},
\end{equation}
which is the exact inverse where $|X|\gg\sqrt\beta$ and has the maximum gain $\max_{|X|}|X|/(|X|^2+\beta) = 1/(2\sqrt\beta) = g$
at $|X| = \sqrt\beta$: no frequency is boosted by more than the chosen \ui{höchste Anhebung} $g$. The filters are blended
to 1 (no processing) outside the band where cancellation can work, $C\leftarrow 1 + a\,b(f)(C-1)$ with the amount $a$ and
a band window $b(f)$ rising by a raised cosine from 40 to \SI{80}{Hz} and falling from 12 to \SI{18}{kHz}. With
\ui{Seitenklang ausgleichen} both filters are also multiplied by the inverse of the near-ear magnitude, smoothed over a
third of an octave and normalised at 500\,Hz--\SI{2}{kHz}, clamped to $\pm\SI{10}{dB}$: it is the same for both ears,
so it commutes with the matrix. Finally both are delayed by 256 samples (\SI{5.3}{ms}) so their non-causal part fits,
transformed to 1024-sample filters with raised-cosine fades at both ends, and run as two FIR convolutions on $m$ and
$s$. The design reports the remaining crosstalk, the ratio of far-ear to near-ear energy between 250\,Hz and
\SI{6}{kHz} computed with the truncated filters against the \emph{unsmoothed} measured $R$:
$e_\mathrm{near/far} = \tfrac12\bigl((1+R)C_s \pm (1-R)C_d\bigr)$.

% ---------------------------------------------------------------------
\subsection{Loudspeakers: the AllRAD decoder (\fn{Speakers.cpp})}
\label{sec:math-allrad}
\fn{buildAllRad} computes a matrix $\mathbf D\in\mathbb R^{S\times K}$ that maps the Ambisonic channels to $S$
loudspeakers, by All-Round Ambisonic Decoding \citep{zotter2012allrad}: decode to a dense, uniform virtual layout,
where Ambisonic decoding is well behaved, and pan each virtual loudspeaker onto the real ones with VBAP
\citep{pulkki1997vbap}, which works for any layout.

\paragraph{Triangulation.} The loudspeaker directions $\mathbf u_s$ are triangulated by their convex hull: a triangle
$(i,j,k)$ belongs to the hull if all other points lie on one side of its plane, tested with the normal
$\mathbf n = (\mathbf u_j-\mathbf u_i)\times(\mathbf u_k-\mathbf u_i)$ (brute force over all triples, fine for up to about
a hundred loudspeakers). If no loudspeaker is above $20^\circ$ (or below $-20^\circ$), an \emph{imaginary} one is added
at the pole; if some directions are still not covered by any triangle (only frontal loudspeakers, say), another
imaginary one is added in the middle of the uncovered part, up to eight times. Imaginary loudspeakers are part of the
triangulation but their signal is discarded, so sound from a direction without loudspeakers is not forced onto the
nearest ones at full level.

\paragraph{VBAP.} For a direction $\mathbf p$ and a triangle with the matrix $\mathbf M = [\mathbf u_i\ \mathbf u_j\
\mathbf u_k]$, the gains are $\mathbf g = \mathbf M^{-1}\mathbf p$; the inverse is stored per triangle from the cross
products, $\mathbf M^{-1} = [\mathbf u_j\times\mathbf u_k,\ \mathbf u_k\times\mathbf u_i,\ \mathbf u_i\times\mathbf u_j]^{\mathsf T}/
\det\mathbf M$. The triangle that contains $\mathbf p$ is the one whose three gains are all non-negative; the gains
are normalised to unit energy, $\mathbf g/\lVert\mathbf g\rVert$.

\paragraph{The virtual decoder with max-$r_E$.} 5000 virtual loudspeakers $\mathbf v_j$ on a Fibonacci lattice
($z_j = 1-2\frac{j+1/2}{5000}$, azimuth $j\cdot\pi(3-\sqrt5)$, the golden angle) are decoded by sampling, each weighted
with the max-$r_E$ window \citep{zotter2019}:
\begin{equation}
  D_{s,k} = \sum_j g_{s}(\mathbf v_j)\; w_{l(k)}\; Y_k(\mathbf v_j),\qquad
  w_l = (2l+1)\,P_l\bigl(\cos\theta_E\bigr),\quad \theta_E = \frac{137.9^\circ}{L+1.51}.
\end{equation}
For a plane wave of SN3D coefficients $Y_k(\mathbf p)$, the signal of virtual loudspeaker $j$ is then
$\sum_l w_l P_l(\cos\gamma_j)$ by the addition theorem (the factor $2l+1$ turns the SN3D product sum into the
full-normalised one), with $\gamma_j$ the angle between $\mathbf v_j$ and $\mathbf p$: a panning lobe around
$\mathbf p$. The weights $P_l(\cos\theta_E)$ shape that lobe so that the energy vector $\mathbf r_E$ is as long as
possible, i.e.\ the sound is as concentrated as the order allows, at the cost of slightly wider main lobes than the
plain (\enquote{basic}) decoder. The Legendre polynomials are evaluated by Bonnet's recurrence
$P_k = \bigl((2k-1)xP_{k-1}-(k-1)P_{k-2}\bigr)/k$.

\paragraph{Normalisation and quality.} The decoder is scaled so that the energy $\sum_s(\mathbf D\mathbf y(\mathbf p))_s^2$
of a plane wave, averaged over the directions where real loudspeakers are (VBAP energy on imaginary ones at most half),
is one: as loud as the stereo and binaural outputs. The page shows the spread of that energy over directions and the
localisation error, the angle between $\mathbf p$ and the energy vector $\mathbf r_E = \sum_s e_s\mathbf u_s$ with
$e_s = (\mathbf D\mathbf y)_s^2$, the direction a listener in the centre perceives at high frequencies.

\paragraph{Distance compensation and switching (\fn{SpeakerRenderer}).} With loudspeakers at different radii $r_s$
the farthest one sets the time: loudspeaker $s$ is delayed by $(r_\mathrm{max}-r_s)/c$ and scaled by $r_s/r_\mathrm{max}$
(the $1/r$ law), so all wavefronts arrive together and equally loud. A new decoder (another layout or order) is
crossfaded linearly over \SI{30}{ms} against the old one, channel by channel, including loudspeakers that only the old
one had.
